\uuid{bCJb}
\exo7id{4166}
\titre{exo7 4166}
\auteur{quercia}
\organisation{exo7}
\datecreate{2010-03-11}
\isIndication{false}
\isCorrection{true}
\chapitre{Fonction de plusieurs variables}
\sousChapitre{Dérivée partielle}
\cpge{mpsi-mp2i.fonctions-deux-variables.b, pcsi.fonctions-deux-variables.b, ptsi.fonctions-deux-variables.b, bcpst1.fonctions-deux-variables.c, mp-mpi.calcul-differentiel-optimisation.a}
\module{Analyse}
\niveau{L2}
\difficulte{}

\contenu{
\question{
Montrer que l'égalité $2e^{x+y} + y - x = 0$ définit $y = \varphi(x)$ au voisinage
    de $(1,-1)$. Calculer $\varphi'(1)$ et $\varphi''(1)$.
}
\reponse{
$\varphi'(1) = -\frac13$, $\varphi''(1) = -\frac8{27}$.
}
}
